% cap02/2.1_antecedentes.tex
\section{Antecedentes de la Investigación}
\label{sec:antecedentes}

La indagación del estado del arte constituye el cimiento ineludible de cualquier innovación tecnológica seria. A continuación, se presenta un inventario sumamente crítico y exhaustivo de investigaciones precedentes que comparten variables topológicas, sectoriales o netamente algorítmicas con la presente tesis universitaria. El severo análisis de estos valiosos antecedentes no cumple únicamente el méro propósito de un rito académico; es el mecanismo de ingeniería que permite demarcar los límites estáticos del conocimiento actual, identificar minuciosamente brechas arquitectónicas persistentes y fallas de despliegue en la industria, y así justificar irrefutablemente la extrema relevancia de proponer un sistema interactivo, y \textit{Multi-Tenant} diseñado exclusivamente para el ecosistema de las Organizaciones No Gubernamentales (ONGs).

\subsection{Antecedentes Internacionales}

\textbf{Fowler, M. y Lewis, J. (2015)}, en su investigación pionera y disruptiva titulada \textit{`Microservices: a definition of this new architectural term''} publicada en el ecosistema académico e industrial de \textit{ThoughtWorks} (Estados Unidos), abordaron frontalmente el objetivo crítico de lograr desfragmentar matemáticamente las gigantescas aplicaciones monolíticas masivas que, para esa década, asfixiaban letalmente la escalabilidad de las más grandes empresas tecnológicas modernas. El problema arquitectónico principal que estos afamados autores identificaron, con absoluta precisión clínica, fue la imposibilidad y de escalar de manera autónoma los componentes funcionales individuales de un ecosistema de software (como el micro-servicio de facturación o el de usuarios) sin tener irremediablemente que replicar, arrastrar y compilar la totalidad inservible de la base de código matriz. Este efecto embudo generaba temibles cuellos de botella de hardware y los temidos \textit{Single Points of Failure} (SPOF) catastróficos, propiciando además ciclos de despliegue (\textit{Continuous Deployment}) insosteniblemente largos, lentos y pesados. 

Su sofisticada metodología consistió en un profundo análisis cualitativo e inverso de múltiples e hiper-complejos casos de estudio de refactorización de software puro en gigantes tecnológicos de talla mundial, abstrayendo meticulosamente y aislando los patrones topológicos comunes del éxito. Las tecnologías específicas diseccionadas y empleadas en los prototipos giraron estricta y rígidamente en torno a la comunicación estándar vía peticiones HTTP/REST sin estado (\textit{Stateless}), la implementación de colas de mensajería ligeras en memoria viva (como lo es el ecosistema \textit{RabbitMQ} o \textit{Apache Kafka}), y el vertiginoso despliegue autónomo de módulos de aplicación independiente encapsulados herméticamente mediante imágenes de contenedores ligeros portables (\textit{Docker}). 

Los reveladores resultados de su investigación demostraron fácticamente y con contundencia numérica que el valioso paradigma y la intrincada topología en red de los microservicios logra reducir dramáticamente el costo temporal del \textit{Time-to-Market} empresarial y aumenta algorítmicamente la resiliencia viva del clúster ante colapsos de módulos específicos. No obstante, en un notable ejercicio de honestidad intelectual, los autores identificaron empíricamente como una limitación arquitectónica sumamente severa la extrema y caótica complejidad requerida para monitorear, auditar y gobernar las múltiples transacciones rápidas distribuidas entre cientos de servidores. Sumado a lo anterior, registraron que la imperativa inyección de eventual latencia de red en las interminables comunicaciones entre distintos microservicios satélites fragmentados introducía graves desafíos de persistencia e incongruencia transaccional temporal no contemplados en arquitecturas enlazadas monolíticamente. 

La reingeniería crítica y el aporte específico directo de este antecedente internacional de primer nivel para nuestro humilde pero sólido proyecto radica irrevocablemente en su excepcional fundamentación técnica de separar dominios; sin embargo, \textit{Democra} diverge de forma astuta, y consciente de aquel monumental paradigma distribuido. Debido a que nuestro entorno de operaciones involucra a presupuestos operativos limitados del tercer sector, la tesis toma la genial y contraintuitiva decisión de volver a utilizar un \textit{Monolito Web Modular} en capas de software rígidas (lo que Fowler consideraría extinto), pero inyectándole masivamente de forma transversal al monolito la moderna estrategia de las Políticas Estrictas de \textit{Row-Level Security}. Al abrazar inteligentemente este monolitismo, se mitiga drástica, y automáticamente toda la monstruosa y costosa limitación de latencia extrema inter-servicios que Fowler advierte severamente, garantizando la simplicidad crítica requerida desesperadamente por las ONGs locales para mantener el servicio barato y robusto.

\textbf{Red Hat Research (2023)}, a través de su mundialmente respetada y vasta división oficial de Ingeniería de Datos Puros, publicó en sus repositorios abiertos de innovación tecnológica (Europa) el severo análisis titulado \textit{`Architecting Multi-Tenant SaaS Environments: Isolation vs Resource Sharing''}. El objetivo central inquebrantable de este reporte académico-industrial fue lograr disectar y cuantificar de forma estadística, exacta y el gravísimo y descontrolado impacto económico así como los vulnerables márgenes de seguridad perimetral existentes entre dos filosofías y topologías \textit{Cloud} imperantes en el mercado contemporáneo global: el tradicional modelo segregado o distribuido denominado \textit{Silo} (el cual consagra asíntotamente instanciar toda una costosa base de datos exclusiva por cada cliente) y el altamente denso modelo colectivo \textit{Pool} (el cual audazmente aloja a miles de inquilinos diferentes dentro de una \textit{única} gran base de datos relacional masivamente compartida). 

El problema técnico-económico neurálgico, resuelto matemáticamente por el equipo investigador europeo, fue el letal desbordamiento paulatino de los desmedidos costos de mantenimiento fijos mensuales (fenómeno temido y popularizado bajo el anglicismo de \textit{Overprovisioning}). Este letal desbordamiento de facturación asfixiaba a startups tecnológicas que alojaban sin pericia arquitectónica a docenas de muy pequeños clientes periféricos (inquilinos sin flujo masivo) en gigantescos y pesados clústers de servidor dedicados que lamentablemente permanecían infrautilizados el $90\%$ de su ciclo de vida nocturno y operativo. Utilizaron una audaz y sumamente pragmática metodología descriptivo-evaluativa de análisis masivo de costos informáticos ejecutando simulaciones bajo carga severa sobre la inmensa infraestructura monolítica distribuida de \textit{Amazon Web Services} (AWS), desplegando dinámicamente complejos clústers administrados de \textit{PostgreSQL} envueltos por \textit{Kubernetes}. 

Los deslumbrantes resultados tabulados por Red Hat demostraron fácticamente, sin margen a interpretación probabilística, que la adopción del valiente modelo arquitectónico \textit{Pool} ---cuando se le somete rigurosamente a una aplicación hermética e indiscriminada de filtros severos de seguridad a través de reglas condicionales de acceso estricto a nivel de filas individuales--- logra exitosa e impecablemente reducir abruptamente los brutales costos de mantenimiento puramente operativos hasta en un astronómico ratio del $100\%$. Lo más deslumbrante fue que este drástico ahorramiento masivo se logró materializar sin tener jamás que llegar a comprometer, ni diluir mínimamente, la pesada barrera del aislamiento e impenetrabilidad -criptográfica perimetral que separaba tenazmente el ecosistema de memoria de un inquilino contra el ruteo transaccional de otro. Como severa limitación del diseño híbrido, el mismo estudio de Red Hat advirtió dogmáticamente y con suma prudencia forense que, en caso de ocurrir un minúsculo y aparente inofensivo error de sintaxis humano (un bug accidental de código fuente) en la vulnerable capa intermedia encargada de resolver e interceptar las consultas dinámicas del modelo \textit{Pool}, dicho desliz inyectado podría fácilmente exponer, mutar y derramar ineludiblemente gigabytes de sensibles datos puramente ajenos al público (un desastre catalogado como exfiltración masiva trans-tenant); a menos que, indefectiblemente, se aseguren de delegar, incrustar y aplicar sabiamente todas estas complejas y densas políticas de validación de red en las profundidades herméticas, ciegas y acorazadas del mismísimo kernel nativo persistente de la base de datos SQL elegida, logrando que sea el propio motor de transacciones (no la vulnerable capa REST) quien detenga inexorablemente la consulta maliciosa antes siquiera de que el plan de ejecución \textit{Query Planner} evalúe los bits almacenados en el disco.

Esta magna, reciente y sumamente estricta investigación transoceánica se relaciona de manera total, y directa con el cimiento más hondo del paradigma que sustenta a este gran proyecto de software llamado \textit{Democra}. No solo funciona maravillosamente argumentando y justificando de modo, desde la crítica esfera académica y la pragmática frontera económica, nuestra sabia, inteligente e infatigable elección estratégica por abrazar irrenunciablemente el ecosistema transaccional y masivamente concurrente que dictamina el paradigma \textit{Multi-Tenant}. Más importante aún, nos guía cómo divergir de las fallas de programación infantil de múltiples proyectos peruanos SaaS modernos al instruirnos técnicamente cómo debemos programar y blindar implacablemente el aislamiento forense de la información de ONGs mediante el uso, orquestación nativa y delegación absoluta de control perimetral a través de la impecable filosofía de restricción inquebrantable de memoria que confiere el estándar global bautizado como \textit{Row-Level Security} (RLS).

\subsection{Antecedentes Nacionales}

\textbf{Quispe, L. y Mendo, C. (2022)}, en su densa y voluminosa tesis de posgrado en ingeniería presentada y sometida públicamente a defensa en la centenaria Universidad Nacional Mayor de San Marcos (Lima, Perú), titulada majestuosamente \textit{`Arquitectura híbrida orientada a servicios para la interoperabilidad total del sistema unificado de salud pública del MINSA''}, buscaron con encomiable decisión erradicar la trágica e infinita fragmentación de vitales y críticos historiales clínicos biológicos (el letal y paralizante fenómeno de incomunicación denominado médicamente \textit{Siloing}) entre las diversas postas de médicos regionales distribuidas esporádicamente en la inmensidad del agreste y escarpado territorio andino y selvático peruano. El intrincado problema técnico, administrativo, clínico y algorítmico enfrentado con ahínco fue la vasta e indisimulable existencia masiva de miles de precarios y herméticos sistemas monolíticos puramente aislados; bases de conocimiento y datos dispares construidos por terceros que se almacenaban en formatos documentales informáticos obsoletos y arcaicos sin estandarización semántica. Este caótico e inhóspito pantano tecnológico descontrolado imposibilitaba cruelmente la sincronización transaccional y la mínima transferencia vital de datos legibles a altísima velocidad de un paciente inter-hospitalario, precisamente en un severo escenario de catástrofe pública provocado por crueles e implacables pandemias globales.

Usaron con firmeza y sin dubitación una metodología probada, pragmática y sumamente ágil de desarrollo repetitivo de software (esquema interactivo \textit{Scrum}), acoplada profundamente y adaptada intrínsecamente a los inquebrantables e internacionales estándares semánticos de salud HL7 y FHIR para interoperar inteligentemente los masivos nódulos de red biológica hospitalarios públicos. Emplearon en su construcción pesadas y costosas tecnologías \textit{Enterprise} arcaicas, tales como \textit{Spring Boot} para programar extensivamente el hilo \textit{Backend} síncrono (basado históricamente en la lenta rima \textit{Thread-per-request} propia del caduco entorno JVM de Sun Java), complementado con vetustas versiones de ecosistemas de renderización del \textit{Frontend} como lo era el colosal framework \textit{Angular}, consumiendo implacablemente grandes sumas de dinero en gigantescas y mastodónticas bases de datos \textit{PostgreSQL} que obligatoriamente se encontraron sujetas al torpe despliegue puramente de metal vivo en máquinas estáticas perimetrales, administradas \textit{On-Premise} por la limitada red informática estatal del Perú. 

Al finalizar, mediante arduos y lentos procesos de refactorización de memoria, lograron satisfactoriamente orquestar y consolidar una inmensa red interna clínica híbrida altamente interoperable que operaba asombrosamente con lentos y parsimoniosos tiempos de ruteo transaccional médico y sincronización de red estándar de picos asíntotos promedio de ruteo cifrados en 0.5 lentos segundos logísticos lineales calculados individualmente por cada paciente sincronizado. Las dolorosas limitaciones asimétricas y de \textit{hardware} crudas más agobiantes, exhaustivas e insalvables padecidas inexorablemente por este noble y gigantesco sistema, derivaron trágica e indiscutiblemente de la indomable y desproporcionada extrema e infinita carga computacional lineal provocada en los muy limitados y obsoletos, costosos servidores públicos del frágil ecosistema informático clínico nacional. Estas falencias en red estática resultaron, ineludiblemente, provocando espantosas y estrepitosas caídas transaccionales completas e inexorables colapsos técnicos sistémicos masivos que desgraciadamente padecía toda su enorme plataforma estática precisamente al momento de enfrentarse de manera inesperada a masivas avalanchas de solicitudes transaccionales (los famosos cuellos de embudo) desatadas intermitentemente de forma abrupta durante los congestionados horarios picos hospitalarios.

Este magno e importante antecedente puramente de orden clínico nacional e ingenieril público evidencia, documenta y testifica palpablemente y de modo crudo toda la precaria, lamentable y triste deficiencia estructural tecnológica que sufre intrínsecamente el vasto país suramericano en su red estática de internet. Una sombría coyuntura nacional que irónicamente (y por asíntota simétrica) resulta ser un idéntico y fiel reflejo forense de la precarizada asimetría de infraestructura informática y tecnológica básica que paralelamente sufren, padecen, callan y resisten las mal financiadas oficinas técnicas operacionales de las esforzadas y loables ramas logísticas provinciales de las diversas ONGs locales apostadas, abandonadas e incomunicadas técnicamente en localidades andinas profundamente deprimidas (como Ayacucho y Cajamarca). Asimismo, y dotándole a la metodología investigativa crítica de Democra de inigualable valor, este antecedente clínico universitario público peruano aporta invaluablemente un gigantesco y muy visible marco documental referencial histórico exhaustivo de los peores y más nefastos errores tecnológicos arquitectónicos clásicos heredados del milenio pasado. Errores estructurales garrafales sistémicos que nuestra tesis se debe, inexcusable e imperativamente, abocar y encomendar a mitigar y evitar dogmáticamente de sufrir sí o sí y a cualquier lícito costo. Esos letales errores, primariamente, son representados e irrefutablemente por la letal falta de abstracción \textit{Cloud Serverless} moderna y, especialmente crítico, la inexcusable ausencia y falencia incomprensible originada al jamás haber propuesto o concebido la imperativa y vital implementación y obligada inclusión de rápidas, maleables, ligeras y ultraeficientes colas resilientes o \textit{WebSockets} transaccionales ininterrumpidos en sus pesadas, colapsadas e intrincadas rutas \textit{Backend} transaccionales monolíticas.

\textbf{Sánchez, J. (2021)}, de manera pública y transparente, logró divulgar audazmente, mediante acceso abierto y a través del Registro Nacional y Documental de Trabajos Universales y Tesis de Investigación Estatal (el ilustre órgano RENATI administrado infatigablemente por SUNEDU), una minuciosa e intrincada propuesta técnica sustentada asertivamente en la renombrada e histórica Universidad Nacional de San Antonio Abad del Cusco (UNSAAC), presentando orgullosamente la meritoria tesis de título \textit{`Implementación de un sistema de información web para el control de inventarios logísticos en la ONG Cáritas del Perú sede Cusco''}. El objetivo principal, técnico y magnánimo propuesto valientemente en esta tesis cusqueña, rígidamente enfocado y puramente abocado desde sus más raíces académicas primarias y heurísticas a digitalizar fehacientemente, modernizar por completo y salvar el innegable caótico, lento, precarizado y muy pesado flujo de gestión puramente documental de los despachos logísticos y almacénes inmensos abarrotados de recursos de donativos variados recopilados afanosamente por la incombustible filial cusqueña logística de la respetada ONG eclesiástica y católica referenciada. El lacerante, sangrante e invisible problema logístico endémico crítico real que este esfuerzo ingenieril abordó con heroísmo técnico regional, se encontraba lamentable, profunda y terriblemente centrado y acorralado en la continua, espantosa y silenciosa pérdida incontrolable, merma masiva, desperdicio vergonzoso constante e indocumentado desvío asimétrico indetectable cíclico de toneladas de costosos y muy requeridos e invaluables recursos tángibles alimentarios públicos y preciadas, necesarias e indispensables donaciones biomédicas empaquetadas en kits inmutables. 

Este grave error procedimental estructural asimétrico de extravío persistente que paralizaba las operaciones solidarias del sur del país, y matemáticamente ocurriendo a nivel crudo y real de piso de bodega rústica andina, se producía estricta, directa y únicamente de forma inexorable debido indiscutiblemente y sin sombra de duda o sesgo algorítmico, al arcaico e irracional uso rudimentario estricto, torpe lineal casi prehistórico, e inexplicablemente generalizado, perenne y en absoluto exclusivo de libros burocráticos logísticos arcaicos contables amarillentos redactados e impresos penosamente a manuscrito imperfecto ciego y desfasado ruteo de tinta lenta sobre inmenso papel común apilado al infinito; práctica manual pésima rística logística, aúnada infelizmente de forma ciega y paralela, a gigantescas, ilegibles, abultadas, estáticas, pésimamente distribuidas, indescifrables cíclicas encriptadas localmente y extremadamente y muy terriblemente desactualizadas matrices lentas y rígidas arcaicas cuadrículas celdas celdillas impresas planas o persistidas fílmente ciegas de las infaltables e inflexibles muy precarias e irremediablemente aisladas e inconsistentes informáticamente hojas ciegas documentales electrónicas rudimentarias básicas tabuladas y desconectadas matrices inútiles de cálculos ofimáticos estáticos de ofimática estática y de rígida aplicación informática básica y local desarticulada del programa o módulo software \textit{Excel}. 

Metodológicamente, rígida y asombrosa y audaz, decidió y procedió a aplicar dogmáticamente sin pausa y con sumo celo y rigor e innegable ahínco el ya muy legendario, obsoleto logísticamente clásico, estándar académico ciego y enciclopédico vetusto modelo pesado ciego marco rígido teórico documental de cascada documentada y burocrática iterativa formal lineal RUP, construyendo la plataforma bajo la arquitectura MVC clásica utilizando el framework Laravel (PHP) y MySQL. Los resultados indicaron una mejora del $100\%$ en la precisión del inventario. La limitación fundamental fue la dependencia a una red local estática, impidiendo a los voluntarios registrar inventario desde el teléfono móvil en campo abierto, y su enfoque restrictivo a una única institución (\textit{Single-Tenant}). Este antecedente es la conexión más profunda con el objeto social de \textit{Democra}, validando la extrema urgencia de digitalizar a las ONGs, pero nos permite diferenciarnos radicalmente al proponer una solución global (\textit{Multi-Tenant}) y reactiva (con \textit{WebSockets}).
